% Section 5 — Results
% Status: full-depth (14pp Wearable AI target). Body carries Table 1 (headline),
% Table 2 (multi-corpus verdicts), Fig 3 (qualitative), Fig 4 (PSM vs MLLM),
% Fig 5 (HDD). Subsections: Setup,
% Single-session (+Table 1, Fig 3), Memory/latency/failure, Cross-corpus
% generalization (+Table 2), Why MLLMs need a prefilter (+Fig 4), HDD persistence
% (+Fig 5). Supplement keeps ablation prose, memory crossover, HLL micro-eval.

\section{Results}
\label{sec:results}

We evaluate \psm{} against brute-force CLIP-cosine retrieval and a frontier MLLM on egocentric look-back QA (retrieval baselines, i.e. sliding-window and uniform-sample CLIP, are in \S\ref{sec:results-ablations}).

\subsection{Setup}
\label{sec:results-setup}

We evaluate on the Nymeria atomic-action benchmark~\cite{ma2024nymeria}, 30 sessions (${\sim}5{,}600$ narrations) on Project Aria Gen~1 headsets~\cite{engel2024aria}; each narration's $[t_\text{start},t_\text{end}]$ interval is retrieval ground truth, a proxy benchmark, as the narrations are third-person action descriptions, not natural look-back questions (no purpose-built outdoor egocentric look-back QA benchmark exists; \S\ref{sec:conclusion}).
Primary results use \texttt{shelby\_arroyo\_act0} (187 narrations, 69\,m bbox), encoder CLIP-L~\cite{clip-l-laion,radford2021clip} (768-d, 1\,fps).
\psm{} uses the \S\ref{sec:method} defaults ($r{=}12$, $R{=}128$); a full-reservoir row at $R{=}1024$ isolates reservoir eviction from the \caponek{} effect, and the baseline is brute-force CLIP-cosine top-5 over the full $N$-frame bank (no spatial structure, no capping).
We report three top-5 metrics: Hit@5 (any exemplar timestamp in a GT interval), exemplar mIoU@5 (IoU of exemplar $\pm1.5$\,s windows vs.\ GT), and bucket mIoU@5 (IoU of each cell's $[t_\text{min},t_\text{max}]$ envelope vs.\ GT).

\subsection{Fixed-budget allocation: which places get remembered}
\label{sec:results-main}

Our primary experiment isolates \emph{memory allocation} from encoder quality and
from total memory. Every policy receives an identical global exemplar budget
$M{=}128$ and returns the same number of candidates ($K{=}5$, global cosine, no
per-cell cap), so differences reflect only \emph{which} frames a policy chooses to
retain. We compare, on the 14 street-scale sessions (5 seeds, CLIP-L, $r{=}12$):
streaming \texttt{global\_reservoir} (Algorithm~R over the whole stream),
\texttt{fifo}, \texttt{hybrid}, \texttt{uniform\_time}, an offline
\texttt{semantic\_kcenter} coreset (farthest-first in cosine space; a
diversity diagnostic, not a streaming competitor), and \texttt{spatial\_priority},
a causal exact-$M$ policy that protects under-represented cells. Sessions are the
unit of analysis (seeds averaged within session, then paired session bootstrap).
Alongside Hit@5/mIoU@5 we report \emph{oracle retention coverage} (any retained
timestamp in a GT interval, before ranking) and pre-registered rare/common-place
strata.

\begin{table*}[t]
\centering
\begin{tabular}{lccccc}
\toprule
Policy ($M{=}128$) & Hit@5 & mIoU@5 & oracle & rare-place Hit & common-place Hit \\
\midrule
\texttt{semantic\_kcenter}$^\dagger$ & \textbf{27.8\%} & \textbf{0.206} & 75.9\% & 28.0\% & \textbf{26.0\%} \\
\texttt{uniform\_time}      & 26.0\% & 0.194 & \textbf{91.8\%} & 30.1\% & 21.8\% \\
\texttt{global\_reservoir}  & 23.9\% & 0.183 & 78.6\% & 26.7\% & 20.6\% \\
\texttt{spatial\_priority}  & 23.6\% & 0.183 & 76.7\% & \textbf{31.1\%} & 17.9\% \\
\texttt{hybrid}             & 22.0\% & 0.168 & 68.1\% & 21.9\% & 20.0\% \\
\texttt{fifo}               & 14.6\% & 0.119 & 44.5\% & 13.0\% & 15.2\% \\
\bottomrule
\end{tabular}
\caption{Budget-matched retention on 14 street-scale sessions ($M{=}128$, $r{=}12$,
5 seeds, CLIP-L), seed-averaged. All policies keep the same number of exemplars and
return the same number of candidates. $^\dagger$\texttt{semantic\_kcenter} is an
\emph{offline} diversity diagnostic (uses the full bank), not a streaming policy.
\texttt{spatial\_priority} is a causal probe with area-growing bookkeeping, not a
deployed method. Diagnostics \texttt{spatial\_balanced}/\texttt{visit\_balanced}
(non-causal) omitted for space.}
\label{tab:fixed-budget}
\end{table*}

The result is a characterization, not a win (Table~\ref{tab:fixed-budget},
Fig.~\ref{fig:fixed-budget}).
\textbf{Spatial allocation changes which places are remembered, not aggregate
retrieval.} Paired against the budget-matched \texttt{global\_reservoir},
\texttt{spatial\_priority} is statistically flat on Hit@5 ($-0.3$\,pp, 95\% CI
$[-2.3,2.0]$) and mIoU ($-0.0009$ $[-0.015,0.014]$), but shifts the composition of
what is recalled: rare/low-exposure places $+4.4$\,pp $[-1.3,10.3]$ at the cost of
common places $-2.7$\,pp $[-5.1,-0.4]$. The rare-place effect is \emph{place-specific}:
correct spatial alignment beats coordinate-permutation nulls on rare-place recall
(perm202 $+8.3$\,pp $[4.2,12.7]$, perm303 $+7.7$\,pp $[1.5,13.9]$ significant;
perm101 $+4.0$\,pp $[-0.3,8.6]$ marginal) and is robust to grid translation and
$30^\circ$ rotation (Hit@5 $\Delta{\approx}0$, all n.s.), so it is not a
partition-balancing or H3-boundary artifact. Crucially, a non-spatial coreset
(\texttt{semantic\_kcenter}) is \emph{best overall} and beats
\texttt{spatial\_priority} on aggregate Hit@5 ($+4.1$\,pp $[1.5,6.7]$): generic
visual diversity, not physical space, explains the strongest budget-matched
retention. At a tighter budget ($M{=}64$) the spatial policy's rare-place and
aggregate edge over the global reservoir grows (a secondary budget interaction),
but it still trails \texttt{uniform\_time}/\texttt{semantic\_kcenter} overall
(\S\ref{sec:results-ablations}). We therefore present \texttt{spatial\_priority} as
an experimental retention probe and report all baselines and controls in the main
comparison.

\begin{figure}[t]
\centering
\includegraphics[width=\linewidth]{figures/fixed_budget.pdf}
\caption{Budget-matched retention (14 street-scale sessions, CLIP-L, seeds
averaged within session then macro-averaged).
\textbf{(A)}~Hit@5 vs.\ logical exemplar memory ($M\times(768{\cdot}4{+}8)$\,B) for
the four foregrounded policies; the preregistered $M{=}128$, $r{=}12$ operating
point is marked. A non-spatial semantic-diversity coreset (offline) is strongest;
the spatial probe only edges the global reservoir at the tightest budget.
\textbf{(B)}~At $M{=}128$, $r{=}12$, the spatial probe's Hit@5 shifts by stratum
relative to the budget-matched global reservoir --- up on rare/low-exposure places,
down on common places --- with paired session-bootstrap 95\% CIs (sessions
resampled). Not a Pareto frontier.}
\label{fig:fixed-budget}
\end{figure}

\subsection{Single-session operating points}
\label{sec:results-single}

Table~\ref{tab:headline} reports the single-session comparison on \texttt{shelby\_arroyo\_act0} (187 narrations).

\begin{table*}[t]
\centering
\begin{tabular}{lcccc}
\toprule
Method & Reservoir $R$ & Hit@5 & exemplar mIoU@5 & bucket mIoU@5 \\
\midrule
Brute-force CLIP-L     & --   & \textbf{13.4\%} (25/187) & 0.074 & 0.074 \\
PSM, \texttt{cap}=1    & 128  & 9.1\% (17/187)           & 0.048 & 0.019 \\
PSM, \texttt{cap}=2    & 128  & 9.6\% (18/187)           & 0.050 & 0.014 \\
PSM, \texttt{cap}=3    & 128  & 10.2\% (19/187)          & 0.057 & 0.014 \\
PSM, \texttt{cap}=$K$  & 128  & \textbf{11.2\%} (21/187) & 0.063 & 0.013 \\
\midrule
PSM, \texttt{cap}=1    & 1024 & 8.0\% (15/187)           & 0.044 & 0.019 \\
PSM, \texttt{cap}=$K$  & 1024 & \textbf{13.4\%} (25/187) & 0.074 & 0.013 \\
\bottomrule
\end{tabular}
\caption{\psm{} vs brute-force CLIP on \texttt{shelby\_arroyo\_act0} (187 narrations), single representative seed.
$R{=}128$ (top block) is the bounded-memory deployment; $R{=}1024$ isolates reservoir eviction from the \caponek{} effect, recovering brute-force exactly at \caponek{}$=K$. Optional MLLM rerank: \S\ref{sec:results-ablations}.}
\label{tab:headline}
\end{table*}

At the deployment point ($R{=}128$, \caponek{}$=K$), \psm{} reaches $84\%$ of brute-force Hit@5 ($11.2\%$ vs $13.4\%$; mIoU $0.063$ vs $0.074$) with session-length-invariant per-cell state; at $R{=}1024$ ($8\times$ the per-cell reservoir) it recovers the \emph{same 25/187} with identical mIoU, a substrate-correctness check whose $2.2$\,pp/$0.011$ gap to $R{=}128$ is the bounded-memory cost. On this single 187-question session the $2.1$\,pp Hit@5 gap sits within one binomial standard error ($95\%$ Wilson CIs $[7.5,16.6]$ and $[9.2,19.0]$ overlap heavily), so we lean on the seed-averaged 30-session sweep (\S\ref{sec:results-ablations}) for the aggregate claim.
Strict spatial diversity (\caponek{}$=1$) costs $5.4$\,pp at $R{=}1024$ ($8.0\%$ vs $13.4\%$; $4.3$\,pp at $R{=}128$) because revisits to a cell collapse into one bucket, the intended behavior for place-distinct queries.\footnote{Table~\ref{tab:headline} is a single representative seed; its one non-monotonicity (\caponek{}$=1$ Hit@5 higher at $R{=}128$, $9.1\%$, than $R{=}1024$, $8.0\%$) is within reservoir-sampling noise (max-cosine need not be the ground-truth frame, so a random subsample can incidentally surface a GT-interval exemplar the full set's top-1 misses); the seed-averaged 30-session sweep (\S\ref{sec:results-ablations}) is monotone in $R$ at every \caponek{}.}
Brute-force itself caps at $13.4\%$: 162/187 narrations are revisit-heavy and visually indistinguishable to CLIP cosine~\cite{weller2025limit}.

Fig.~\ref{fig:f7-qualitative} shows representative retrievals (\psm{}'s top-5 candidates with the Gemini-selected frame) spanning two hits (one a Gemini Hit@1 rescue) and an honest revisit-heavy miss.

\begin{figure}[t]
\centering
\includegraphics[width=\linewidth]{figures/fig3_qualitative.pdf}
\caption{Qualitative look-back retrievals on \texttt{shelby\_arroyo\_act0}.
Each row is one query: the narration text, GT interval, and \psm{}/Gemini verdict (left), then \psm{}'s top-5 candidates ranked by CLIP cosine (similarity and timestamp per frame). Green/red borders mark whether a candidate falls in the GT interval; the Gemini-selected frame is highlighted.
Rows~1--2 are hits (row~2 a Gemini Hit@1 rescue over \psm{}'s ranking); row~3 is an honest revisit-heavy miss, where recurring near-identical frames defeat cosine retrieval (\S\ref{sec:limitations}).}
\label{fig:f7-qualitative}
\end{figure}

\subsection{Ablations, retrieval baselines, and reranking}
\label{sec:results-ablations}

\paragraph{Operating-point Pareto and non-spatial baselines.}
As \caponek{} sweeps from 1 to $K$, Hit@5 rises and bucket mIoU@5 falls monotonically: the place-diversity/embedding-precision trade the knob is designed to expose. Against non-spatial baselines on the headline session, sliding-window CLIP at $\ell{=}10$\,s edges \psm{} \caponek{}$=K$ by $+1.6$\,pp Hit@5 ($15.0\%$ vs $13.4\%$) by averaging out per-frame variance, but requires the full $N$-frame bank; uniform-sample CLIP forms the floor at $4.3$--$8.6\%$. Tightening the reservoir from $R{=}1024$ to $128$ costs $2.2$\,pp Hit@5 at \caponek{}$=K$, while H3 resolution barely moves Hit@5 (${\pm}0.5$\,pp across $r\in\{10,11,12,13\}$) but moves bucket mIoU monotonically ($0.004\to0.019$).

\paragraph{30-session monotonicity.}
The single-seed headline (Table~\ref{tab:headline}) generalizes to a seed-averaged sweep over all 30 Nymeria sessions. Mean Hit@5 climbs monotonically in \caponek{} at both reservoirs ($2.99\%\to4.81\%\to6.11\%\to7.55\%$ at $R{=}128$ and $3.26\%\to5.69\%\to7.08\%\to8.92\%$ at $R{=}1024$), with the $R{=}1024$ endpoint matching brute-force exactly and the sweep monotone in $R$ at every \caponek{} (e.g.\ $3.26\%$ vs $2.99\%$ at \caponek{}$=1$). This confirms the lone single-seed non-monotonicity in Table~\ref{tab:headline} is reservoir noise. Reservoir eviction thus costs $1.37$\,pp Hit@5 in the 30-session aggregate at \caponek{}$=K$; sliding-window CLIP exceeds both by $+1.0$\,pp.

\paragraph{Optional MLLM rerank.}
An MLLM reranker only \emph{reorders} \psm{}'s fixed top-5, so every @$k$ metric
(Hit@5, mIoU@5) is permutation-invariant and cannot change by construction; an
earlier draft's ``mIoU@5 $0.074{\to}0.101$'' rerank gain was an artifact of a
wider scoring tolerance on the rerank run and is withdrawn. The only rank-sensitive
effect is at the top: reranking moves the MLLM-selected candidate to rank~1, and its
top-1 IoU is small and mixed across sessions (it helps visually-rich queries,
e.g.\ Hit@1 $10\%{\to}30\%$ on \texttt{seq008\_running\_001}, but hurts
action-narration sessions such as \texttt{seq003} $-3.3$\,pp and Nymeria
$-1.6$\,pp). We therefore report reranking only through top-1 metrics, do not
promote it to a default, and note \psm{}'s results hold without any MLLM in the loop.

\subsection{Memory, latency, and failure modes}
\label{sec:results-memory-latency}

\psm{}'s per-cell state is session-length-invariant ($R{=}128\times768$-d$\times4$\,B $\approx0.38$\,MiB plus a 60-slot $p{=}10$ HLL ring $\approx0.06$\,MiB), whereas brute-force grows linearly in $N$.
On the short Nymeria sessions ($N\in[748,1325]$, ${\sim}22$\,min) brute-force is still cheaper. Whether \psm{} is smaller than a dense bank depends on ingest rate and on how the per-cell HLL ring is charged: counting the full $C{=}60$ ring at every visited cell, \psm{}'s \emph{exemplar} term is sublinear but its total state can \emph{exceed} the dense bank at low fps on wide-area corpora (\S\ref{sec:results-hdd}). We therefore report exemplar and ring bytes separately and treat HDD memory as modeled logical state, not a headline bounded-memory win.
Retrieval is sub-millisecond: brute-force dot-product over the 1133-frame bank is 67\,$\mu$s median; the paired in-process \psm{} benchmark is 456/153\,$\mu$s at \caponek{}$=1$/$K$ (host-CPU).
We caution that a previously reported ``${\sim}20$\,MiB, sub-2\,ms'' Galaxy S22
figure was measured at a reduced benchmark configuration ($d{=}128$, $R{=}4$,
$C{=}12$), \emph{not} the CLIP-L deployment point; it is withdrawn as a deployment
number. At the actual deployment configuration ($d{=}768$, $R{=}128$, $C{=}60$) our
host benchmark (20k observations over 500 cells) uses ${\sim}96$\,MiB RSS and answers
an \emph{uncentered} semantic query in ${\sim}13$\,ms (the global scan of
\S\ref{sec:method}; a spatially-centered k-ring query is far cheaper), with the
per-exemplar state ${\sim}192\times$ larger than the reduced config. On-device
figures at the true configuration require an S22 rerun and are deferred; the CLIP
encoder, not the index, remains the dominant on-device cost
(\S\ref{sec:limitations}).
The \caponek{} latency ordering flips with density: on the dense Nymeria bench \caponek{}$=1$ is slower (its $K$ distinct cells each force a separate time-envelope merge, the dominant cost), while on a sparse workload \caponek{}$=K$ is slower (its larger candidate sort dominates); cosine scoring is identical across settings.
Of the 162 shared brute-force/\psm{} misses, most are revisit-heavy narrations with recurring near-identical frames; the 25 groundable questions reference temporally-distinct events, so the shared $13.4\%$ ceiling reflects single-frame CLIP-cosine matching (sliding-window averaging and SigLIP-2 exceed it; \S\ref{sec:results-ablations}, \S\ref{sec:results-multi-corpus}), not method quality (\S\ref{sec:limitations}).

\subsection{Cross-corpus, cross-encoder generalization}
\label{sec:results-multi-corpus}

The spatial-axis result generalizes beyond the headline session. To rule out a
single-corpus / single-encoder artifact, we replicate it across \textbf{14
street-scale sessions spanning 3 sensor stacks} (SLOPER4D~\cite{dai2023sloper4d}:
head-mounted LiDAR + DJI Action2; Nymeria: Project Aria + MPS SLAM;
LookOut~\cite{pan2025lookout}: Project Aria across 5 Bay Area locations)
\textbf{and 3 encoder families} (CLIP-L 768-d, CLIP-bigG 1280-d, SigLIP-2-large
1024-d~\cite{tschannen2025siglip2}). For each (session, encoder) the acceptance
criterion is Hit@5 non-decreasing across $r\in\{10,11,12\}$ with $\geq4$\,pp lift
from $r{=}10$ to $r{=}12$ (conservative against binomial noise at 28--30
questions/session). We pre-register \textbf{CLIP-L} (our primary encoder) as the
verdict basis to avoid post-hoc encoder selection: \textbf{9/14 sessions pass}
under CLIP-L (Table~\ref{tab:multi-corpus-h3}). Reporting instead the
\emph{any-encoder disjunction} (a session passes if \emph{any} of the three
encoders satisfies the criterion) gives 13/14, but this is an optimistic upper
bound that selects the successful encoder after the fact; we include it only for
completeness. Per-encoder, SigLIP-2 passes 11/14, bigG 10/14, CLIP-L 9/14; the
lone all-encoder failure, \texttt{BurlingameDT4\_feb5}, is a visually-homogeneous
trajectory (\S\ref{sec:limitations}). We also caution that these SLOPER4D/LookOut
questions are generated from target frames (\S\ref{sec:limitations}), so this is a
spatial-resolution sensitivity check, not held-out look-back QA. No existing egocentric benchmark combines
outdoor spatial trajectories with temporally-grounded text, so we assemble a
three-corpus scaffold with per-corpus extractors while stopping short of claiming
a released benchmark.

\begin{table*}[t]
\centering
\setlength{\tabcolsep}{4pt}
\small
\begin{tabular}{llccccc}
\toprule
& & \multicolumn{3}{c}{Per-encoder PASS} & & \\
\cmidrule(lr){3-5}
Corpus & Session & \texttt{clipL} & \texttt{bigG} & \texttt{siglip2L} & Any-enc & \#PASS \\
\midrule
LookOut  & \texttt{Mainquad\_jan10}              & \cmark & \xmark & \cmark & \cmark & 2/3 \\
LookOut  & \texttt{Sanmateopark\_garage\_jan11}  & \xmark & \xmark & \cmark$\dagger$ & \cmark & 1/3 \\
LookOut  & \texttt{Fostersquare1\_jan16}         & \xmark & \cmark & \cmark & \cmark & 2/3 \\
LookOut  & \texttt{BurlingameDT5\_feb5}          & \xmark & \cmark & \cmark & \cmark & 2/3 \\
LookOut  & \texttt{SanmateoDT2\_Jan12}           & \cmark & \cmark & \cmark & \cmark & 3/3 \\
LookOut  & \texttt{Gates\_to\_mainquad\_jan10}   & \cmark & \cmark & \cmark & \cmark & 3/3 \\
LookOut  & \texttt{Huang\_Gates\_jan10}          & \cmark & \cmark & \cmark & \cmark & 3/3 \\
LookOut  & \texttt{BurlingameDT4\_feb5}          & \xmark & \xmark & \xmark & \xmark & 0/3 \\
LookOut  & \texttt{SSC3\_jan17}                  & \cmark & \cmark & \cmark & \cmark & 3/3 \\
LookOut  & \texttt{Hillsdale6\_jan14}            & \xmark & \xmark & \cmark$\dagger$ & \cmark & 1/3 \\
\midrule
SLOPER4D & \texttt{seq003\_street\_002}          & \cmark & \cmark & \cmark & \cmark & 3/3 \\
SLOPER4D & \texttt{seq008\_running\_001}         & \cmark & \cmark & \cmark & \cmark & 3/3 \\
SLOPER4D & \texttt{seq009\_running\_002}         & \cmark & \cmark & \xmark & \cmark & 2/3 \\
\midrule
Nymeria  & \texttt{shelby\_arroyo\_act0}         & \cmark & \cmark & \xmark & \cmark & 2/3 \\
\midrule
\multicolumn{2}{r}{\textbf{Total PASS}}
                  & \textbf{9/14} & \textbf{10/14} & \textbf{11/14}
                  & \textbf{13/14} & --- \\
\bottomrule
\end{tabular}
\caption{Acceptance verdicts on the multi-corpus, multi-encoder spatial-axis sweep (5 seeds $\times$ 5 H3 resolutions $\times$ 3 encoders per session).
\emph{Criterion}: a session passes if an encoder's Hit@5 is non-decreasing across $r{\in}\{10,11,12\}$ with a $\geq4$\,pp $r10{\to}r12$ lift.
\textbf{Pre-registered CLIP-L is the verdict basis: 9/14 pass.} ``Any-enc'' is the post-hoc disjunction over all three encoders (13/14), an optimistic upper bound reported for completeness; $\dagger$~marks SigLIP rescues where neither CLIP-family encoder passes. \texttt{BurlingameDT4} fails all encoders (visually homogeneous). Questions here are target-frame-generated (\S\ref{sec:limitations}), so this is a resolution-sensitivity check, not held-out QA.}
\label{tab:multi-corpus-h3}
\end{table*}

\subsection{Query-aware retrieval vs.\ uniform frame sampling}
\label{sec:results-vanilla-mllm}

On these same 14 sessions we compare query-aware retrieval against a frontier MLLM
over \emph{uniformly-sampled} frames. We run Gemini 3.1 Pro at $K{=}8$ uniform
frames per question (Fig.~\ref{fig:f2-psm-vs-mllm}). We are careful about what this
measures: Gemini's \emph{actual chosen-frame} hit rate is only ${\sim}5\%$ (mean
over sessions), because with $K{=}8$ uniform frames on 100--640\,s sessions the
gap to a ${\sim}3$\,s GT interval means the right frame is rarely even in the
candidate set. An \emph{any-of-$K$ coverage} score (a hit if \emph{any} of the 8
uniform frames lands in GT, independent of Gemini's pick) is an upper bound of
${\sim}8\%$ and reflects uniform temporal coverage, not the MLLM's reasoning.
\psm{}'s query-guided top-5 exceeds both on nearly every session. The honest
conclusion is narrow: \textbf{query-aware retrieval beats uniform frame sampling};
this does \emph{not} show that MLLMs require \psm{} (a strong MLLM given the right
candidates could do well), only that uniform sampling is a poor way to feed one.

\begin{figure}[t]
\centering
\includegraphics[width=\linewidth]{figures/fig4_psm_vs_mllm.pdf}
\caption{Query-aware \psm{} \texttt{clipL@r12} top-5 (blue) vs.\ $K{=}8$ uniform-frame sampling on the 14 street-scale sessions. Red is any-of-$K$ \emph{coverage} (a hit if \emph{any} of the 8 uniform frames lands in GT) --- an upper bound reflecting temporal coverage, not the MLLM's reasoning; Gemini's \emph{actual} chosen-frame hit rate is lower still (${\sim}5\%$ mean). Query-aware retrieval beats uniform sampling; this does not show MLLMs require \psm{} (\S\ref{sec:results-vanilla-mllm}).}
\label{fig:f2-psm-vs-mllm}
\end{figure}

\subsection{Multi-session persistence at metropolitan scale (Honda HDD)}
\label{sec:results-hdd}

The results so far are single-pass, so they cannot exercise the multi-session
persistence and revisit-cardinality properties (\S\ref{sec:intro}) that motivate
a bounded index. We close this gap on the Honda Research Institute Driving
Dataset (HDD)~\cite{ramanishka2018hdd}: 104\,h of San Francisco Bay Area driving
with RTK GPS and a front camera (our extraction covers 132 drives, 102.9\,h with
valid GPS), which over 8 months genuinely re-drives the same streets. HDD ships
no look-back QA labels, so these results are \emph{systems} and
\emph{self-supervised}, not retrieval-accuracy (\S\ref{sec:limitations}).
Driving is not wearable capture, but it supplies what no egocentric look-back
corpus yet offers (months of GPS-anchored revisits to the same places). We use it
only to characterize \psm{}'s per-cell memory model against a real multi-session
revisit structure; the revisit property turns on structure, not locomotion mode
(cf.\ the GPS-anchored wearable walk in Fig.~\ref{fig:f8-psm-viz}).

\begin{figure}[t]
\centering
\includegraphics[width=0.72\linewidth]{figures/f5_hdd_memory.pdf}
% TODO(figure): regenerate fig5 to a single panel (A only); panels B/C are
% withdrawn as PSM-engine results (offline Python substitutes; see text).
\caption{HDD memory model (132 drives, 102.9\,h). \psm{}'s modeled per-cell
\emph{logical} state (exemplar term) vs.\ a dense per-frame bank over cumulative
driving hours, from the RTK trajectory (no embeddings). The exemplar term is
sublinear in frame count, but charging the full $C{=}60$ HLL ring to every
visited cell makes \psm{}'s total state exceed the bank at low fps (see text);
we plot exemplar and ring terms separately. Earlier cardinality/cross-session-AUC
panels are withdrawn as engine results (offline substitutes;
\S\ref{sec:results-hdd}).}
\label{fig:hdd-combined}
\end{figure}

\paragraph{Revisit structure and a corrected memory model (Fig.~\ref{fig:hdd-combined}A).}
Binning every drive to H3 $r{=}10$ (${\sim}66$\,m), \textbf{47.0\%} of
(drive, cell) traversal events fall in cells revisited on $\geq2$ distinct days
(pre-registered bar $\geq30\%$), across 5{,}163 revisited cells spanning a
\textbf{median of 106 days} (max 221). This is a genuine multi-session revisit
structure. We use it to model \psm{}'s \emph{logical} per-cell state from the RTK
trajectory (no embeddings): the \emph{exemplar} term $\sum_\text{cells}\min(\text{frames},R)$
grows sublinearly in frame count, $1.4$/$2.3$/$4.4$/$7.2\times$ below a dense bank
at 1/5/15/30\,fps. However, charging the full $C{=}60$ HLL ring
($60\,\mathrm{KiB}$) to \emph{every} visited cell --- most seen once --- adds
${\sim}1.4$\,GB over 23{,}298 cells, so \psm{}'s \emph{total} state can
\emph{exceed} the dense bank at 1\,fps ($2.27$ vs $1.14$\,GB) and the advantage is
weaker throughout ($0.5$/$1.5$/$3.2$/$5.5\times$). We therefore report exemplar and
ring bytes separately and treat this as \emph{modeled logical state}, not a
bounded-memory guarantee: total state grows with the cell count, which itself
grows over the 8 months (only 20\% of cells seen by the halfway mark). This
per-cell independent allocation --- a full quota to every cell --- is precisely
what a global budget would avoid.

\paragraph{Engine correctness, and why we demote the corpus cardinality/retrieval claims.}
The HDD embeddings needed to run per-cell cardinality and cross-session retrieval
\emph{through the C engine} at corpus scale are not available (only a single
sanity drive was extracted); an earlier draft's ``$8{,}112$ unique observations''
accrual and ``$0.853$ cross-session AUC'' were computed by \emph{offline Python
substitutes} (a BLAKE2b all-time accumulator, and a dense $\mathbf{e}\!\cdot\!\mathbf{q}$
scan over all frames), not by \psm{}'s reservoir/Murmur-HLL engine. Because those
numbers do not exercise the released engine at scale and rest on features we cannot
re-run, \textbf{we withdraw them as PSM-engine results} and defer engine-backed
multi-session HDD cardinality/retrieval to future work. What we \emph{can} verify
is an implementation check: on the one available drive the engine's per-cell HLL
matches the Python re-implementation to a median $0.05\%$ (and both to the true
distinct count to $0.00\%$ with a non-aging window), and a Nymeria micro-eval puts
per-session HLL error at median $1.0\%$. We also note (\S\ref{sec:method}) that
for continuous CLIP embeddings this cardinality counts distinct \emph{frames}, not
distinct events or people, absent an object/event canonicalization stage.
